\section{Introduction}
Detectors such as YOLO~\cite{ultralytics2023} run in real time on single-board
computers, but a deployment is not a benchmark. A Raspberry Pi~5 in a closed
enclosure reaches its firmware thermal limit within minutes at full load; once it
throttles, frame rate drops and jitters. Memory pressure, a weak supply
(under-voltage), a frozen camera or a hung process produce other silent failures.
Operators typically pick one model for the worst case (wasting capacity) or the
best case (failing in the field).

\shade{} treats the model configuration as a controllable resource. Our
contributions are:
\begin{enumerate}
\item A \textbf{control plane} for commodity boards: a degradation ladder,
  hysteresis-guarded transitions, atomic background hot-swap with a measured
  stall, a heartbeat supervisor with last-known-good restart, and a SQLite audit
  trail of decisions and reasons (Sec.~\ref{sec:design}).
\item A \textbf{content-aware scheduler} that gates detection on scene change
  and propagates boxes through skipped frames by per-box phase correlation,
  without exceeding the ladder's detection budget (Sec.~\ref{sec:sched}).
\item An \textbf{evaluation method}: a closed-loop thermal simulator that
  drives the real decision engine, with ablations and a protocol to fit it to a
  device (Sec.~\ref{sec:sim}).
\item \textbf{Measurements and negative results} on when adaptation pays,
  including a learned thermal-forecast policy that did not beat tuned reactive
  control in simulation (Secs.~\ref{sec:results}, \ref{sec:sim}).
\end{enumerate}
We do not claim a new model-switching or DVFS algorithm; the aim is a small,
reproducible, auditable system and an honest account of its benefits.
Hardware results are \TBD{pending Pi~5 runs E1--E4, E7--E9}; simulation and
offline results are labelled as such.
